\section{Related Work}
\label{sec:related}
% =====================================================================

We position PG-DSL against five adjacent literature lanes.

\paragraph{MCP description-channel defences.}
MCPShield~\cite{mcpshield} is the closest neighbour on the
\emph{(pre-invocation $\times$ tool-description-verification)} axis. It
uses a three-stage LLM-as-judge pipeline (cognitive probing, isolated
projection, periodic reasoning) and reports a $95.30\%$ defence rate
over digital MCP threats. PG-DSL differs along four axes
simultaneously: (i) physics-simulation verifier vs.\ LLM judge, (ii)
CPS substrate vs.\ digital tools, (iii) theorem-bearing analysis vs.\
empirical-only, and (iv) automatic NL-to-formal lifting vs.\ informal
prompt judgement. Crucially, MCPShield's evaluation does \emph{not}
include any attack with a physical effect, leaving the
CPS-physical attack class structurally outside its tested coverage.
``Securing the Model Context Protocol''~\cite{secure_mcp_jamshidi}
provides taxonomies and detection heuristics for general MCP threats
but does not address the description--physics gap on CPS substrates.
Recent measurement work~\cite{mcp_misleading} quantifies real-world
MCP server description--code mismatch (the threat
PG-DSL targets) but does not propose a defence;
PG-DSL is complementary as a defensive companion to that finding.

\paragraph{Contract-grounded tool execution.}
ToolGate~\cite{toolgate} formalises each tool as a Hoare-style contract
(precondition, postcondition) and verifies state evolution against the
contracts at runtime. ToolGate's contracts are \emph{provided}; PG-DSL
\emph{lifts} them from natural-language descriptions. ToolGate verifies
symbolic state at runtime; PG-DSL verifies physical state in a digital
twin at admission. The two are complementary at the framing level but
address different layers of the agent stack.

\paragraph{Classical CPS intrusion detection.}
INVARLLM~\cite{invarllm} extracts physical invariants offline using an
LLM and detects runtime violations on iTrust SWaT and WADI with
$100\%$ precision. Sparse strong observability under joint
sensor/actuator attacks~\cite{showkatbakhsh_sparse} provides the
formal foundation for runtime state-anomaly detection. Both are
runtime, state-conditional defences; neither addresses
description-channel attacks. They are PG-DSL's natural composition
partners (Theorem~\ref{thm:t3}).

\paragraph{NL-to-formal lifting.}
Req2LTL~\cite{req2ltl} demonstrates LLM-based lifting of
natural-language software requirements into LTL with $88.4\%$ semantic
accuracy. PG-DSL specialises this to a CPS-grounded grammar
(\S\ref{sec:method-grammar}); T2 builds on Req2LTL's framework but
proves soundness conditions for the restricted class. Req2LTL's
accuracy gap ($1 - 0.884 = 0.116$) provides the executor's draft
anchor for $\varepsilon_L$.

\paragraph{Theoretical foundations and adversarial framing.}
The trust-authorisation mismatch SoK~\cite{trust_auth_sok} establishes
the impossibility of static agent verification.
``Gaming the Judge''~\cite{gaming_judge} establishes the
manipulability of LLM-judge verification. Together these motivate
\emph{physics-grounded tool verification} as the defensible escape.
MSB~\cite{msb_mcp_bench} and ``When MCP Servers
Attack''~\cite{when_mcp_attack} catalogue attack classes against
MCP-controlled agents; the prompt-injection survey
~\cite{prompt_inj_survey} surveys
description-channel attacks broadly. PG-DSL's threat model
(\S\ref{sec:threat}) inherits the description-channel attack
categorisation from this lane and operationalises a
physics-substrate-grounded defence.

% =====================================================================
